\section{Empirical Strategy}\label{coll}

We bring five detection approaches to the conviction-anchored benchmark: a close-margin regression discontinuity (Section~\ref{sec:rdd}), classical bid-distribution screens (Section~\ref{sec:classical}), network screens built from worker-flow and co-bidding data (Section~\ref{sec:network}), a specialization placebo (Section~\ref{sec:placebo}), and---crucially---a temporal-conditioning test that distinguishes prospective from retrospective detection power (Section~\ref{sec:temporal_method}). Every screen is evaluated by its area under the ROC curve (Section~\ref{sec:auc}). A simple model in Section~\ref{sec:model} organizes predictions about where classical screens should fail and where post-enforcement labor reallocation should concentrate.

\subsection{Close-margin regression discontinuity}\label{sec:rdd}

Following \citet{kawai2023using}, we focus on auctions decided by a hair. Define the margin of victory in auction $a$ as
\[
MV_a \;=\; \frac{b^{\text{runner-up}}_a - b^{\text{win}}_a}{b^{\text{win}}_a}\,.
\]
When $MV_a$ is close to zero, the winner barely underbid the nearest rival. Under genuine competition, narrowly winning and narrowly losing firms should look alike on pre-determined characteristics; a sharp discontinuity at the cutoff signals something non-competitive.

We estimate the jump using local-polynomial regression with a data-driven bandwidth. Two outcomes carry interpretive weight. \emph{Incumbency}---whether the firm won the previous auction in the same market---tests for bid rotation: a jump would indicate that designated winners are systematically past winners. \emph{Last bidder}---whether the firm submitted the final bid---tests for strategic timing. A density test on the running variable checks whether the margin itself is being manipulated.

\subsection{Classical bid-distribution screens}\label{sec:classical}

We compute four screens on firm-best bids (one bid per firm per auction, keeping the lowest when multiple rounds exist):

\begin{enumerate}
\item \textbf{Coefficient of variation} ($CV_a$): standard deviation divided by mean. Bid rotation predicts lower $CV$ because cover bids cluster near the designated winner \citep{porter1993detecting, abrantesmetz2006variance}.

\item \textbf{Bid spread} ($S_a$): $(\max - \min)/\text{mean}$. Like $CV$, compressed under cover bidding.

\item \textbf{Margin of victory} ($MV_a$): the designated winner should undercut cover bids by just enough to win, yielding a small margin.

\item \textbf{Number of bidders} ($N_a$): entry deterrence or market division should reduce the count \citep{marshall2012economics}.
\end{enumerate}

For each, we compare cartel-tainted auctions to clean controls on the close-margin benchmark ($\lvert MV \rvert \le 0.10$), reporting the Welch $t$-statistic, Cohen's $d$, the Wilcoxon rank-sum $p$-value, and the AUC.

\subsection{Network-based screens}\label{sec:network}

We step outside the bid data and build three screens from administrative sources.

\noindent\textbf{Worker-flow screens.} For each (winner, runner-up) pair we measure labor connectivity from the full 2009--2017 RAIS panel:

\begin{itemize}
\item \textbf{Shared workers} ($W_{ij}$): the count of distinct PIS numbers appearing at both firms during the nine-year window.

\item \textbf{Jaccard similarity} ($J_{ij} = W_{ij} / (N_i + N_j - W_{ij})$): normalizes the count by firm size, so that large employers are not mechanically flagged.
\end{itemize}

The logic is straightforward. If cartel coordination in thin labor markets runs through personnel---managers carrying pricing information, technicians rotating between member firms---then the employment records should show it \citep{stoyanov2012worker, serafinelli2019good}. Alternatively, if the signal is retrospective---reflecting labor reallocation \emph{after} the cartel period ends---it would still appear in the full-panel measure but would vanish under temporal conditioning. We test both possibilities.

\noindent\textbf{Co-bidding screen.} For each pair we count the auctions in which both firms submitted bids ($C_{ij}$). This provides a network channel built entirely within the procurement data.

\subsection{Detection evaluation: AUC}\label{sec:auc}

We use the area under the ROC curve to compare screens. The AUC is simply the probability that a randomly chosen cartel auction scores more ``cartel-like'' than a randomly chosen clean auction. For screens where lower is more suspicious ($CV$, $S$, $MV$, $N$), we negate scores so that AUC $> 0.5$ always means detection in the predicted direction. We compute the AUC via the Mann--Whitney $U$ statistic with average-rank tie-handling:
\[
\text{AUC} \;=\; \frac{U_1}{n_1 \times n_0}\,,
\]
where $n_1$ and $n_0$ are the cartel and clean sample sizes. All screens are evaluated on the same unified close-margin benchmark, so comparisons are apples-to-apples. Appendix~\ref{app:bandwidth} shows that results are stable across bandwidths from $0.02$ to $0.50$.

\subsection{Temporal conditioning and decomposition}\label{sec:temporal_method}

Network screens built from administrative panel data face a fundamental identification challenge: the measure uses information from the \emph{entire} panel, including periods after the auction being screened. If cartel detection or enforcement changes the network---for instance, by triggering labor reallocation from convicted firms to competitors---then the full-panel AUC conflates prospective detection power with retrospective correlation. We implement two tests.

\emph{Temporal conditioning.} For each auction at year $t$, we recompute the worker-flow and co-bidding screens using only data from years strictly before $t$. If the AUC survives, the screen captures information available at the time of the auction. If it collapses, the signal is retrospective. We report a cumulative backward window $[2009, t-1]$ and a rolling window $[\max(2009, t-3), t-1]$.

\emph{Temporal decomposition.} For each shared worker between a cartel firm and a rival, we classify the inter-firm link by when it was first observable relative to the cartel's conduct period: \textbf{pre-conduct} (both firms employed the worker before the start year), \textbf{during-conduct} (the link appeared during $[\text{start}, \text{end}]$), or \textbf{post-conduct} (the link appeared after the end year). We then use each window's shared-worker count as a separate screen and evaluate it against the benchmark. This decomposition identifies \emph{when} the signal arises---and whether it reflects coordination, ongoing cartel activity, or enforcement aftermath.

\subsection{Specialization placebo}\label{sec:placebo}

Elevated worker flows between a cartel firm and its rival could reflect several forces: pre-existing personnel connections, the natural circulation of a thin labor pool, or post-conduct labor reallocation. To separate market-level circulation from the cartel-specific component, we build a three-way comparison. The \textbf{cartel} sample consists of auction-level (winner, runner-up) pairs with at least one cartel-active firm ($N \approx 2{,}300$). The \textbf{placebo} consists of non-cartel pairs bidding on the same BEC item code in the same year. The \textbf{clean} baseline is every other non-cartel pair.

If the entire worker-flow signal were market specialization, the placebo should match the cartel sample. Any excess in the cartel sample over the placebo is a cartel-specific residual---though, as we show in the results, this residual is driven primarily by post-conduct displacement rather than by coordination during the cartel period.

\subsection{A framework for screen heterogeneity}\label{sec:model}

Why should screen effectiveness depend on the market? And why should post-enforcement labor reallocation concentrate in thin markets? Consider $n \ge 2$ firms operating in a product market with specialized-labor pool of thickness $L$.

The first question has a standard answer. Classical bid-distribution screens exploit the compression that bid rotation imposes on the bid distribution. In thick markets (large $L$, many bidders), rotation is the natural coordination device and the compression is detectable. In thin markets (small $L$, few bidders), the bid distribution is inherently noisy---few firms generate high variance even under collusion---and phantom-bidding or cover-bid mechanisms may inflate rather than compress dispersion. Classical screens should therefore work in thick markets and fail in thin ones.

The second question is about the labor-market aftermath of enforcement. When a cartel firm contracts---losing contracts, facing fines, undergoing investigation---it sheds workers. In thick labor markets, those workers disperse across many potential employers; the per-firm flow is dilute and hard to detect. In thin markets, the same workers have few alternative employers, and the flows concentrate in the small set of firms that operate in the same specialized product market. The worker-flow screen, evaluated on the full panel, captures this concentration---but it is a concentration of \emph{displacement}, not of \emph{coordination}.

Four testable predictions follow:

\begin{enumerate}
\item \textbf{Coordination-technology gradient.} Classical screens exploit bid-distribution compression---a signature of bid rigging. Coordination that operates outside the auction (e.g., hub-and-spoke arrangements where a central intermediary sets prices upstream) does not compress the bid distribution and should be invisible to classical screens. Classical AUCs should therefore be high for bid rigging and near chance for hub-and-spoke cartels.

\item \textbf{Thickness gradient for classical screens.} In thin markets, classical AUCs should be near or below chance even for bid rigging, because few bidders generate high variance mechanically. In thick markets, they should recover. This prediction is independent of the temporal-conditioning critique and uses contemporaneous bid data.

\item \textbf{Thickness gradient for post-conduct reallocation.} Full-panel worker-flow AUCs should be highest in thin markets---not because coordination runs through personnel, but because displacement concentrates where the labor pool is shallowest.

\item \textbf{Temporal asymmetry.} If the worker-flow signal reflects coordination, it should be present in pre-conduct flows. If it reflects enforcement aftermath, it should be concentrated in post-conduct flows. The temporal decomposition directly tests this.
\end{enumerate}

The framework is deliberately spare. It abstracts from repeated-game dynamics, leniency, and heterogeneity across auction formats. We offer it not as a structure to be estimated but as a lens for organizing the results and for specifying \emph{when} each family of screens should and should not carry signal.
